\documentclass[12pt]{article}
\usepackage[utf8]{inputenc}
\usepackage{amsmath, amssymb} % Math packages
\usepackage{geometry}
\usepackage{hyperref} % Clickable links
\usepackage{graphicx} % Images

\geometry{a4paper, margin=1in}

\title{logica}
\author{Matteo}
\date{\today}

\begin{document}

\maketitle

\section{deduzione naturale}

una regola senza premesse si dice assioma

due tipi di passi di interferenza
\begin{itemize}
\item regole di introduzione di un connettivo, modi in cui concludere direttamente una formula
\item regole di eliminazione, modi in cui utilizzare direttamente un'ipotesi
\end{itemize}

ogni passo di inferenza puo essere letto sia top-down(dalla conclusione alle premesse) che bottom-up(dalle premesse alla conclusione)

una regolsa e invertibile quando dimostra anche ogni premessa

è consigliabile dimostrare con metodo top-down, stando attenti a usare regole invertibili o usando regole non invertibili e acertandosi che la conclusione sia ancora dimostrabile.

\textbf{and introduzione}:
bottom-up: se F1 e F2 allora F1 and F2
top-down: per dimostrare F1 and F2 devo dimostrare sia  F1 che F2

\textbf{and elimininazione}:

bottom-up: se F1 and F2 e se ipotizzando F1 e F2 concludo F3 allora F3

top-down: per dimostrare F3 data l'ipotesi F1 and F2 è sufficente dimostrare F3 sotto le ipotesi f1 e F2

\textbf{or introduzione}

bottom-up: se F1 o F2 vale allora vale anche F1 or F2

top-down: per dimostrare F1 or F2 è sufficente dimostrare F1 o F2

\textbf{or eliminazione}
bottom-up: se vale F1 or F2 e F3 vale sia quando vale F1 che quando vale F2 , allora f3 vale.

top-down: per dimostrare qualunque cosa sapendo F1 or F2, è sufficente procedere per casi, dimosttrando la stessa cosa assumendo prima che f1 valga e poi che valga F2.

\textbf{bottom introduzione}: NESSUNA

\textbf{bottom eliminazione}: 
bottom up: dal falso segue qualunque cosa

top-down: per dimostrare qualunque cosa posso ridurmi a dimostrare un assurdo.

\textbf{=> introduzione}:
bottom-up: se ipotizzando F1 dimostro F2 allora F1=>F2

top-down:  per dimostrare F1 =>F2 basta assumere F1 e dimostrare F2

\textbf{=> eliminazione}:
bottom-up: se F1 e F1=>F2 allora necessariamente F2

top-down: per dimostrare F2 devo trovare un F1 che valga e tale per cui F1=>F2

\textbf{negazione introduzione}: 
bottom-up: se ipotizzando F1 dimostro l'assurdo allora non f1

top-down: per dimostrare non F1 basta assumere F1 e dimostrare l'assurdo.

\textbf{negazione introduzione}: 
bottom-up: è assurdo avere sia non F1 che f1

top-down: per dimostrare l'assurdo basta dimostrare qualcosa e il suo contrario.

quando ci si trova a dimostrare bottom, da quel momento tutte le regole sono invertibili.

\textbf{ ragionamento per assurdo}: 
bottom-up: assumiamo per assurdo non F.... assurdo! quindi F

top-down: per dimostrare F procediamo per assurdo assumendo non F e dimostrando bottom.

\section{semantica classica}

\textbf{connotazione}: una parte di frase sintatticamente corretta alla quale attribuire un significato

\textbf{denotazione }: il significato attribuito a una connotazione, preso da un dominio di interpretazione.

\textbf{dominio di interpretazione}: insieme di possibili significati

\textbf{sintassi}: descrizione dell'insieme di tutte le connotazioni

\textbf{funzione semantica}: funzione che associa a ogni connotazione una denotazione in un dominio di interpretazione fissato..


grande differenza tra semantica classica e intuizionistica:
classica permette di usare RAA e terzo escluso, concludenedo un'ipotesi dimostrando l'assurdo della negazione.
intuizionistica non permette di usare queste due tecniche, permette solo di dimostrare la negazione di non A a partire da A e dimostrando bottom usando non introduzione.

lo stesso esercizio puo essere svolto con semantica classica o intuizionistica.

tutto questo perchè la semantica classica ha come fondamento che se A non e vero allora non a, l'intuizionistica no.

\section{primo ordine}

la \textbf{logica proposizionale}: è quella logica le cui connotazioni possono denotare valori di  verità.

la \textbf{logica del prim''ordine} ha \textbf{formule}: connotazioni che denotano valori di verità

e \textbf{termini}: connotazioni che denotano elementi del dominio del discorso A

La logica si chiama “del primo ordine” perché le variabili (tipo x) rappresentano solo elementi concreti di un insieme. no funzioni o proprieta

\textbf{binder}: lega una variabile in uno scope

i quantificatori sono un tipo di binder.

i binder catturano le varibili rendendole non piu libere. esistono anche variabili libere.

alfa-convertibilità = cambiare nome alle variabili legate

puoi cambiare il nome di una varibile legata senza problemi ma non una variabile libera

\section{deduzione naturale del prim'ordine}

\textbf{$\forall$ introduzione}:
se dimostro P per un elemento completamente arbitrario(y, che deve essere nuovo e non usato nell'ipotesi), allora P vale per tutti gli elementi

Lettura bottom-up: se P vale per una y qualunque, allora per
ogni x vale P

Lettura top-down: per dimostrare $\forall$x.P è sufficiente scegliere
una y sulla quale non sappiamo nulla e poi dimostrare P[y/x].
sia A che B. Ovviamente è sufficiente scegliere una y ancora
mai usata.

\textbf{$\forall$ eliminazione}:
se vaale per tutti allora vale per un caso specifico

Lettura bottom-up: Lettura bottom-up: se P vale per tutti gli x, allora vale in
particolare per t.

Lettura top-down:per dimostrare P[t/x] è sufficiente
dimostrare il teorema generalizzato $\forall$x.P.


\textbf{$\exists$ introduzione}:
se vale per un caso, allora esiste qualcuno per cui vale

è fortemente non invertibile

Lettura bottom-up: se P vale per t allora esiste un x per cui P
vale.

Lettura top-down: per dimostrare $\exists$x.P bisogna scegliere un t
per il quale P[t/x] valga e dimostrarlo

\textbf{$\exists$ eliminazione}:
“Prendo un elemento qualunque che soddisfa P (senza sapere chi è), e vedo cosa posso dedurre”

parzialmente invertibile

Lettura bottom-up: se $\exists$x.P e se dimostro C sotto l’ipotesi che
P valga per un generico y, allora C vale.

Lettura top-down: per dimostrare un qualche C sotto l’ipotesi
$\exists$x.P è sufficiente dimostrare C assumendo P per una qualche
variabile generica y.


quando devi dimostrare un ipotesi, puoi trovarti a dover usare induzione(ricorsione).
cio consiste nel assumere un ipotesi induttiva, e usarla per dimostrare tutti i casi.

slides: 4:deduzione nat(L1)-9:deduz nat prim'ordine(L5)-10:ricorsione induzione(L2+L3) piu imortanti.poi 3:relations functions(L4),7:prim'ordine(capireL5). 2:prove-8:proprieta connettivi per supporto. su 1 solo assioma intersezione, unione e insieme potenza.





\end{document}